%\title{University of Bristol Thesis Template}
\RequirePackage[l2tabu]{nag}		% Warns for incorrect (obsolete) LaTeX usage
%
%
% File: memoirthesis.tex
% Author: Victor Brena
% Description: Contains the thesis template using memoir class,
% which is mainly based on book class but permits better control of 
% chapter styles for example. This template is an adaptation and 
% modification of Oscar's.
% 
% Memoir is a flexible class for typesetting poetry, fiction, h
% non-fiction and mathematical works as books, reports, articles or
% manuscripts. CTAN repository is found at:
% http://www.ctan.org/tex-archive/macros/latex/contrib/memoir/
%
%
% UoB guidelines for thesis presentation were found at:
% http://www.bris.ac.uk/esu/pg/pgrcop11-12topic.pdf#page=49
%
% UoB guidelines:
%
% The dissertation must be printed on A4 white paper. Paper up to A3 may be used
% for maps, plans, diagrams and illustrative material. Pages (apart from the
% preliminary pages) should normally be double-sided.
%
% Memoir class loads useful packages by default (see manual).
\documentclass[a4paper,11pt,leqno,oneside]{memoir} %add 'draft' to turn draft option on (see below)
%
%
% Adding metadata:
\usepackage[usenames,dvipsnames,svgnames,table]{xcolor}
\usepackage{calligra}
\usepackage{tikz}
\usetikzlibrary{arrows,arrows.meta,positioning,calc,patterns,decorations.pathreplacing,fit,shapes,snakes,
		       automata,backgrounds,
		       petri,topaths}%geometric/algebraic description.
%\usepackage[version=0.96]{pgf}			%PGF/TikZ is a tandem of languages for producing vector graphics from a 

\usepackage{tcolorbox}

\newtcolorbox{bidentidad}[2][]{
  enhanced,
  skin=enhancedlast jigsaw,
  attach boxed title to top left={xshift=-4mm,yshift=-0.5mm},
  fonttitle=\bfseries\sffamily,
  colbacktitle=darkgray,
  colframe=orange,
  interior style={
    top color=lightblue!3,
    bottom color=blue!3
  },
  boxed title style={
    empty,
    arc=0pt,
    outer arc=0pt,
    boxrule=0pt
  },
  underlay boxed title={
    \fill[darkgray] 
      (title.north west) -- 
      (title.north east) -- 
      +(\tcboxedtitleheight-1mm,-\tcboxedtitleheight+1mm) -- 
      ([xshift=4mm,yshift=0.5mm]frame.north east) -- 
      +(0mm,-1mm) -- 
      (title.south west) -- cycle;
    \fill[darkgray] 
      ([yshift=-0.5mm]frame.north west) -- 
      +(-0.4,0) -- 
      +(0,-0.3) -- cycle;
    \fill[darkgray] 
      ([yshift=-0.5mm]frame.north east) -- 
      +(0,-0.3) -- 
      +(0.4,0) -- cycle; 
  },
}


\newtcolorbox{bidentidad2}[2][]{
  enhanced,
  skin=enhancedlast jigsaw,
  attach boxed title to top left={xshift=-4mm,yshift=-0.5mm},
  fonttitle=\bfseries\sffamily,
  colbacktitle=darkgray,
  colframe=darkgray,
  interior style={
    top color=blue!3,
    bottom color=red!3
  },
  boxed title style={
    empty,
    arc=0pt,
    outer arc=0pt,
    boxrule=0pt
  },
  underlay boxed title={
    \fill[darkgray] 
      (title.north west) -- 
      (title.north east) -- 
      +(\tcboxedtitleheight-1mm,-\tcboxedtitleheight+1mm) -- 
      ([xshift=4mm,yshift=0.5mm]frame.north east) -- 
      +(0mm,-1mm) -- 
      (title.south west) -- cycle;
    \fill[darkgray] 
      ([yshift=-0.5mm]frame.north west) -- 
      +(-0.4,0) -- 
      +(0,-0.3) -- cycle;
    \fill[darkgray] 
      ([yshift=-0.5mm]frame.north east) -- 
      +(0,-0.3) -- 
      +(0.4,0) -- cycle; 
  },
title = More generally;
}

\usepackage{enumitem}
\usepackage{pifont}
\tcbuselibrary{skins}
\usepackage{psvectorian}
\renewcommand*{\psvectorianDefaultColor}{Blue}%
\usepackage{datetime}
\usepackage{ifpdf}
\ifpdf
\pdfinfo{
   /Author (Peterson Yuhala)
   /Title (MII ENSPY)
   /Keywords (Xen; )
   /CreationDate (D:\pdfdate)
}
\fi
% When draft option is on. 
\ifdraftdoc 
	\usepackage{draftwatermark}				%Sets watermarks up.
	\SetWatermarkScale{0.3}
	\SetWatermarkText{\bf Draft: \today}
\fi
%
% Declare figure/table as a subfloat.
\newsubfloat{figure}
\newsubfloat{table}
\usepackage{booktabs}
\usepackage{siunitx}

\usetikzlibrary{chains}
%\usepackage{background}
% \backgroundsetup{
% scale=1,
% opacity=1,
% angle=0,
% color=black,
% contents={%
% \begin{tikzpicture}[every node/.style={inner sep=0pt}]
% \node[anchor=north west](CNW)
% at (current page.north west) {\pgfornament[width=1.75cm]{61}};
% \node[anchor=north east](CNE)
% at (current page.north east) {\pgfornament[width=1.75cm,symmetry=v]{61}};
% \node[anchor=south west](CSW)
% at (current page.south west) {\pgfornament[width=1.75cm,symmetry=h]{61}};
% \node[anchor=south east](CSE)
% at (current page.south east) {\pgfornament[width=1.75cm,symmetry=c]{61}};
% \end{tikzpicture}%
%   }
% }

\usepackage{multirow}
\usepackage{algorithm}
\usepackage[noend]{algpseudocode}
\makeatletter
\def\BState{\State\hskip-\ALG@thistlm}
\makeatother

% Better page layout for A4 paper, see memoir manual.
\settrimmedsize{297mm}{210mm}{*}
\setlength{\trimtop}{0pt} 
\setlength{\trimedge}{\stockwidth} 
\addtolength{\trimedge}{-\paperwidth} 
\settypeblocksize{634pt}{448.13pt}{*} 
\setulmargins{4cm}{*}{*} 
\setlrmargins{*}{*}{1} 
\setmarginnotes{17pt}{51pt}{\onelineskip} 
\setheadfoot{1.3\onelineskip}{2\onelineskip} 
\setheaderspaces{*}{2\onelineskip}{*} 
\checkandfixthelayout
%
\frenchspacing
% Font with math support: New Century Schoolbook
\usepackage{fouriernc}
\usepackage[nohints]{minitoc}
\usepackage{fancybox}
\usepackage{makeidx}




\definecolor{marron}{RGB}{60,30,10}
\definecolor{orange}{RGB}{241, 110, 0}
\definecolor{darkblue}{RGB}{0,0,80}
\definecolor{blue}{RGB}{0, 0, 128}
\definecolor{lightblue}{RGB}{80,80,80}
\definecolor{darkgreen}{RGB}{0,80,0}
\definecolor{darkgray}{RGB}{25,25,112}
\definecolor{cbordeau}{RGB}{248,217,210}
\definecolor{cblue}{RGB}{245,247,249}
\definecolor{darkred}{RGB}{80,0,0}
\definecolor{midnightblue}{RGB}{25, 25, 112}
\definecolor{shadecolor}{rgb}{0.97,0.97,0.97}
\usepackage{fourier-orns}
\usepackage[T1]{fontenc}


\makeatletter
\def\headrule{{\color{darkgray}\raisebox{-2.1pt}[10pt][10pt]{\leafright} \hrulefill \raisebox{-2.1pt}[10pt][10pt]{~~~\decofourleft \decotwo\decofourright~~~} \hrulefill \raisebox{-2.1pt}[10pt][10pt]{ \leafleft}}}
\makeatother

\usepackage[
final,
stretch=10,
protrusion=true,
tracking=true,
spacing=on,
kerning=on,
expansion=true]{microtype}

\newcommand{\ornamento}{\vspace{-2em} \noindent \textcolor{blue}{\hrulefill~ \raisebox{-2.5pt}[10pt][10pt]{\leafright \decofourleft \decothreeleft  \aldineright \decotwo \floweroneleft \decoone   \floweroneright \decotwo \aldineleft\decothreeright \decofourright \leafleft} ~  \hrulefill}}
\newcommand{\ornpar}{\noindent \textcolor{blue}{ \raisebox{-1.9pt}[10pt][10pt]{\leafright} \hrulefill \raisebox{-1.9pt}[10pt][10pt]{\leafright \decofourleft \decothreeleft  \aldineright \decotwo \floweroneleft \decoone}}}
\newcommand{\ornimpar}{\textcolor{blue}{\raisebox{-1.9pt}[10pt][10pt]{\decoone \floweroneright \decotwo \aldineleft \decothreeright \decofourright \leafleft} \hrulefill \raisebox{-1.9pt}[10pt][10pt]{\leafleft}}}
%
% UoB guidelines:
%
% Text should be in double or 1.5 line spacing, and font size should be
% chosen to ensure clarity and legibility for the main text and for any
% quotations and footnotes. Margins should allow for eventual hard binding.
%
% Note: This is automatically set by memoir class. Nevertheless \OnehalfSpacing 
% enables double spacing but leaves single spaced for captions for instance. 
\OnehalfSpacing 
%
% Sets numbering division level
\setsecnumdepth{subsection} 
\maxsecnumdepth{subsubsection}
%
% Chapter style (taken and slightly modified from Lars Madsen Memoir Chapter 
% Styles document
\usepackage{calc,soul,fourier}
\makeatletter 
\newlength\dlf@normtxtw 
\setlength\dlf@normtxtw{\textwidth} 
\newsavebox{\feline@chapter} 
\newcommand\feline@chapter@marker[1][4cm]{%
	\sbox\feline@chapter{% 
		\resizebox{!}{#1}{\fboxsep=1pt%
			\colorbox{blue}{\color{white}\thechapter}% 
		}}%
		\rotatebox{90}{% 
			\resizebox{%
				\heightof{\usebox{\feline@chapter}}+\depthof{\usebox{\feline@chapter}}}% 
			{!}{\scshape\so\@chapapp}}\quad%
		\raisebox{\depthof{\usebox{\feline@chapter}}}{\usebox{\feline@chapter}}%
} 
\newcommand\feline@chm[1][4cm]{%
	\sbox\feline@chapter{\feline@chapter@marker[#1]}% 
	\makebox[0pt][c]{% aka \rlap
		\makebox[1cm][r]{\usebox\feline@chapter}%
	}}
\makechapterstyle{daleifmodif}{
	\renewcommand\chapnamefont{\normalfont\Large\scshape\raggedleft\so} 
	\renewcommand\chaptitlefont{\normalfont\Large\bfseries\scshape} 
	\renewcommand\chapternamenum{} \renewcommand\printchaptername{} 
	\renewcommand\printchapternum{\null\hfill\feline@chm[2.5cm]\par} 
	\renewcommand\afterchapternum{\par\vskip\midchapskip} 
	\renewcommand\printchaptertitle[1]{\color{blue}\chaptitlefont\raggedleft ##1\par}
} 
\makeatother 

%\renewcommand\refname{\textbf{Reference}}
\chapterstyle{daleifmodif}
%
% UoB guidelines:
%
% The pages should be numbered consecutively at the bottom centre of the
% page.
\renewcommand{\sectionmark}[1]{\markright{\thesection~- ~#1}}
\renewcommand{\chaptermark}[1]{\markboth{Chapitre~\thechapter~-~ #1}{}}
% style sur les numéros de page (avec nom)
\makepagestyle{myvf} 
\makeoddfoot{myvf}{}{\ornimpar \\ \large \hfill \sffamily\bf \textcolor{blue}{ \thepage}
}{} 
\makeevenfoot{myvf}{}{\ornpar   \\ \large  \sffamily\bf \textcolor{blue}{\thepage ~~~ \reflectbox{\leafNE}}  \hfill}{} 
%\makeheadrule{myvf}{\textwidth}{\normalrulethickness} 
\makeevenhead{myvf}{\leftmark}{\ornamento}{\rightmark} 
\makeoddhead{myvf}{\leftmark}{\ornamento}{\rightmark}
\pagestyle{myvf}


%
% Oscar's command (it works):
% Fills blank pages until next odd-numbered page. Used to emulate single-sided
% frontmatter. This will work for title, abstract and declaration. Though the
% contents sections will each start on an odd-numbered page they will
% spill over onto the even-numbered pages if extending beyond one page
% (hopefully, this is ok).
\newcommand{\clearemptydoublepage}{\newpage{\thispagestyle{empty}\cleardoublepage}}
%
%
% Creates indexes for Table of Contents, List of Figures, List of Tables and Index
\makeindex
%\printglossaries 
%below creates a list of abbreviations. \gls and related
% commands are then used throughout the text, so that latex can automatically
% keep track of which abbreviations have already been defined in the text.
%
% The import command enables each chapter tex file to use relative paths when
% accessing supplementary files. For example, to include
% chapters/brewing/images/figure1.png from chapters/brewing/brewing.tex we can
% use
% \includegraphics{images/figure1}
% instead of
% \includegraphics{chapters/brewing/images/figure1}
\usepackage{import}

% Add other packages needed for chapters here. For example:
\usepackage{lipsum}					%Needed to create dummy text
\usepackage{subcaption}
\usepackage{amsfonts} 					%Calls Amer. Math. Soc. (AMS) fonts
\usepackage[centertags]{amsmath}			%Writes maths centred down
\usepackage{stmaryrd}					%New AMS symbols
\usepackage{pgfplots}
\pgfplotsset{yticklabel style={text width=3em,align=right}}

\usepackage{amssymb}					%Calls AMS symbols
\usepackage{amsthm}					%Calls AMS theorem environment
\usepackage{newlfont}					%Helpful package for fonts and symbols
\usepackage{layouts}					%Layout diagrams
\usepackage{graphicx}					%Calls figure environment
\usepackage{longtable,rotating}			%Long tab environments including rotation. 
%\usepackage[latin1,applemac]{inputenc}			%Needed to encode non-english characters 
\usepackage[utf8]{inputenc}
									%directly for mac
\usepackage{color}
\usepackage{colortbl}					%Makes coloured tables
\usepackage{wasysym}					%More math symbols
\usepackage{mathrsfs}					%Even more math symbols
\usepackage{float}						%Helps to place figures, tables, etc. 
\usepackage{verbatim}					%Permits pre-formated text insertion
\usepackage{upgreek }					%Calls other kind of greek alphabet
\usepackage{latexsym}					%Extra symbols
\usepackage[numbers,sort&compress]{natbib}

\usepackage[breaklinks]{hyperref}

%\usepackage[style=authoryear,maxcitenames=2,uniquelist=false,sorting=nyt,backend=bibtex]{biblatex}
%\usepackage[backend=biber]{biblatex}
%\addbibresource{biblio.bib}
%\usepackage[
%backend=biber,
%style=numeric,
%sorting=ynt
%]{biblatex}
%\addbibresource{biblio.bib}

		%Calls bibliography commands 
\usepackage{url}						%Supports url commands
\usepackage{etex}						%eTeXÕs extended support for counters
\usepackage{fixltx2e}					%Eliminates some in felicities of the 
									%original LaTeX kernel
%\usepackage[none]{hyphenat}
%\usepackage[spanish,english,french]{babel}		%For languages characters and hyphenation
\usepackage{memhfixc}					%Must be used on memoir document 
									%class after hyperref
\usepackage{hyperref}
\hypersetup{pdfborder=0 0 0}

\usepackage{color}                    				%Creates coloured text and background
\usepackage[numbers,sort&compress]{natbib}
\usepackage{footnote}					%For footnotes
\usepackage{afterpage}
\usepackage[object=vectorian]{pgfornament}
\usetikzlibrary{shapes.geometric,calc}
\definecolor{fondpaille}{cmyk}{0,0,0.1,0}
\usepackage{pagecolor}
\usepackage{microtype}					%Makes pdf look better.
\usepackage{rotfloat}					%For rotating and float environments as tables, 
									%figures, etc. 
\usepackage{alltt}						%LaTeX commands are not disabled in 
									%verbatim-like environment

				%To use diverse features from tikz	
         
\tcbset{
    Baystyle/.style={
        sharp corners,
        enhanced,
        boxrule=6pt,
        colframe=OliveGreen,
        height=600pt,
        width=500pt,
        borderline={8pt}{-11pt}{},
    }
}


%							
%Reduce widows  (the last line of a paragraph at the start of a page) and orphans 
% (the first line of paragraph at the end of a page)
\widowpenalty=1000
\clubpenalty=1000
%
% New command definitions for my thesis
%
\newcommand{\keywords}[1]{\par\noindent{\small{\bf Keywords:} #1}} %Defines keywords small section
\newcommand{\parcial}[2]{\frac{\partial#1}{\partial#2}}                             %Defines a partial operator
\newcommand{\vectorr}[1]{\mathbf{#1}}                                                        %Defines a bold vector
\newcommand{\vecol}[2]{\left(                                                                         %Defines a column vector
	\begin{array}{c} 
		\displaystyle#1 \\
		\displaystyle#2
	\end{array}\right)}
\newcommand{\mados}[4]{\left(                                                                       %Defines a 2x2 matrix
	\begin{array}{cc}
		\displaystyle#1 &\displaystyle #2 \\
		\displaystyle#3 & \displaystyle#4
	\end{array}\right)}
\newcommand{\pgftextcircled}[1]{                                                                    %Defines encircled text
    \setbox0=\hbox{#1}%
    \dimen0\wd0%
    \divide\dimen0 by 2%
    \begin{tikzpicture}[baseline=(a.base)]%
        \useasboundingbox (-\the\dimen0,0pt) rectangle (\the\dimen0,1pt);
        \node[circle,draw,outer sep=0pt,inner sep=0.1ex] (a) {#1};
    \end{tikzpicture}
}
\newcommand{\ra}[1]{\renewcommand{\arraystretch}{#1}}
\newcommand{\myrowcolour}{\rowcolor[gray]{0.925}}

\newcommand{\range}[1]{\textnormal{range }#1}                                             %Defines range operator
\newcommand{\innerp}[2]{\left\langle#1,#2\right\rangle}                                 %Defines inner product
\newcommand{\prom}[1]{\left\langle#1\right\rangle}                                         %Defines average operator
\newcommand{\tra}[1]{\textnormal{tra} \: #1}                                                       %Defines trace operator
\newcommand{\sign}[1]{\textnormal{sign\,}#1}                                                   %Defines sign operator
\newcommand{\sech}[1]{\textnormal{sech} #1}                                                  %Defines sech
\newcommand{\diag}[1]{\textnormal{diag} #1}                                                    %Defines diag operator
\newcommand{\arcsech}[1]{\textnormal{arcsech} #1}                                       %Defines arcsech
\newcommand{\arctanh}[1]{\textnormal{arctanh} #1}                                         %Defines arctanh
%Change tombstone symbol

\newcommand{\blackged}{\hfill$\blacksquare$}
\newcommand{\whiteged}{\hfill$\square$}
\newcounter{proofcount}
\renewenvironment{proof}[1][\proofname.]{\par
 \ifnum \theproofcount>0 \pushQED{\whiteged} \else \pushQED{\blackged} \fi%
 \refstepcounter{proofcount}
 \normalfont 
 \trivlist
 \item[\hskip\labelsep
       \itshape
   {\bf\em #1}]\ignorespaces
}{%
 \addtocounter{proofcount}{-1}
 \popQED\endtrivlist
}
%
%
% New definition of square root:
% it renames \sqrt as \oldsqrt
\let\oldsqrt\sqrt
% it defines the new \sqrt in terms of the old one
\def\sqrt{\mathpalette\DHLhksqrt}
\def\DHLhksqrt#1#2{%
\setbox0=\hbox{$#1\oldsqrt{#2\,}$}\dimen0=\ht0
\advance\dimen0-0.2\ht0
\setbox2=\hbox{\vrule height\ht0 depth -\dimen0}%
{\box0\lower0.4pt\box2}}
%
% My caption style
\newcommand{\mycaption}[2][\@empty]{
	\captionnamefont{\scshape} 
	\changecaptionwidth
	\captionwidth{0.9\linewidth}
	\captiondelim{.\:} 
	\indentcaption{0.75cm}
	\captionstyle[\centering]{}
	\setlength{\belowcaptionskip}{10pt}
	\ifx \@empty#1 \caption{#2}\else \caption[#1]{#2}
}
\newcommand*\circled[1]{\tikz[baseline]{\node[shape=circle,draw,inner sep=2pt] (char) {#1};}}

%
% My subcaption style
\newcommand{\mysubcaption}[2][\@empty]{
	\subcaptionsize{\small}
	\hangsubcaption
	\subcaptionlabelfont{\rmfamily}
	\sidecapstyle{\raggedright}
	\setlength{\belowcaptionskip}{10pt}
	\ifx \@empty#1 \subcaption{#2}\else \subcaption[#1]{#2}
}
%
\newcommand{\sectionlinetwo}[2]{%
  \nointerlineskip \vspace{.5\baselineskip}\hspace{\fill}
  {\color{#1}
    \resizebox{0.5\linewidth}{2ex}
    {{%
    {\begin{tikzpicture}
    \node  (C) at (0,0) {};
    \node (D) at (9,0) {};
    \path (C) to [ornament=#2] (D);
    \end{tikzpicture}}}}}%
    \hspace{\fill}
    \par\nointerlineskip \vspace{.5\baselineskip}
  }
%An initial of the very first character of the content
\usepackage{Carrickc,ArtNouvc,Kramer,Konanur,Rothdn}
\usepackage{lettrine}
%\renewcommand\LettrineFontHook{\Rothdnfamily}
\newcommand{\initial}[1]{%
	\lettrine[lines=3,lhang=0.33,nindent=0em]{
		\color{blue}
     		{\textsc{#1}}}{}}
     		
\newcommand{\corner}[1]{%
  \begin{tikzpicture}[color=darkgray,remember picture, overlay]
    \node[anchor=north west] at (current page.north west){%
      \pgfornament[width=2cm]{#1}};
    \node[anchor=north east] at (current page.north east){%
      \pgfornament[width=2cm,symmetry=v]{#1}};
    \node[anchor=south west] at (current page.south west){%
      \pgfornament[width=2cm,symmetry=h]{#1}};
    \node[anchor=south east] at (current page.south east){%
      \pgfornament[width=2cm,symmetry=c]{#1}};
  \end{tikzpicture}%
}
\newcommand{\cornerplus}[2]{%
  \begin{tikzpicture}[color=darkgray,remember picture, overlay]
    \node[anchor=north west] at (current page.north west){%
      \pgfornament[width=2cm]{#1}};
    \node[anchor=north east] at (current page.north east){%
      \pgfornament[width=2cm,symmetry=v]{#1}};
    \node[anchor=south west] at (current page.south west){%
      \pgfornament[width=2cm,symmetry=h]{#1}};
    \node[anchor=south east] at (current page.south east){%
      \pgfornament[width=2cm,symmetry=c]{#1}};
    \node[anchor=north] at (current page.north){%
      \pgfornament[width=6.5cm,symmetry=h]{#2}};
    \node[anchor=south] at (current page.south){%
      \pgfornament[width=6.5cm]{#2}};
  \end{tikzpicture}%
}
\newcommand{\pt}[1]{%
  \begin{tikzpicture}[color=darkgray,remember picture, overlay]
    \node[anchor=north] at (current page.north){%
      \pgfornament[symmetry=h]{#1}};
    \node[anchor=south] at (current page.south){%
      \pgfornament[symmetry=h]{#1}};
  \end{tikzpicture}%
}
%
% Theorem styles used in my thesis
%
\AtBeginDocument{\renewcommand{\bibname}{Références}}
\AtBeginDocument{\renewcommand{\refname}{Références}}

\theoremstyle{plain}
\newtheorem{theo}{Theorem}[chapter]
\theoremstyle{plain}
\newtheorem{prop}{Proposition}[chapter]
\theoremstyle{plain}
\theoremstyle{definition}
\newtheorem{dfn}{Definition}[chapter]
\theoremstyle{plain}
\newtheorem{lema}{Lemma}[chapter]
\theoremstyle{plain}
\newtheorem{cor}{Corollary}[chapter]
\theoremstyle{plain}
\newtheorem{resu}{Result}[chapter]
%
% Hyphenation for some words
%
\hyphenation{res-pec-tively}
\hyphenation{mono-ti-ca-lly}
\hyphenation{hypo-the-sis}
\hyphenation{para-me-ters}
\hyphenation{sol-va-bi-li-ty}
%
%
\usepackage[acronym, toc]{glossaries}
\makeglossaries

%%%%%%%%%%%%%%%%%%%%%%%%Abrreviations%%%%%%%%%%%%%%%%%


%bord de la première de couverture

\newcommand{\bord} {
	\begin{tikzpicture}[remember picture, overlay, start chain, node distance=-2mm]
	\node (nworn) [shift={(1mm,-1mm)}, anchor=north west, on chain ] at (current page.north west){\pgfornament[color=blue,width=10mm]{61}};
	\foreach \i in {1,...,7}
	{%\node [on chain] {\pgfornament[width=10mm]{62}};
		\node [on chain] {\pgfornament[color=blue,width=25.5mm]{71}};}
	%\node (neorn) [on chain] {\pgfornament[width=8mm]{62}};
	{\node [on chain] {\pgfornament[color=blue,width=10mm]{62}};}
	\foreach \i in {1,...,12}
	{\node [continue chain=going below, on chain] {\pgfornament[color=blue,width=10.5mm,symmetry=h]{62}};
		\node [continue chain=going below, on chain] {\pgfornament[color=blue,width=10mm]{62}};}
	\node (seorn) [on chain] {\pgfornament[color=blue,width=10mm,symmetry=h]{62}};
	\foreach \i in {1,...,7}
	{\node [continue chain=going left, on chain]{\pgfornament[color=blue,width=25.5mm]{71}};}
	\node (sworn) [on chain] {\pgfornament[color=blue,width=10mm,symmetry=h]{61}};
	\foreach \i in {1,...,12}
	{\node [continue chain=going above, on chain] {\pgfornament[color=blue,width=10.5mm]{61}};
		\node [continue chain=going above, on chain] {\pgfornament[color=blue,width=10mm,symmetry=h]{61}};}
	\end{tikzpicture}
}



\setlength{\parindent}{0.6cm}

\usepackage{chngcntr}
\usepackage[section]{placeins}  
\usepackage{titletoc}

\usepackage[french,english]{babel}


\usepackage[french]{babel}
\addto\captionsfrench{%
  \renewcommand{\contentsname}{Table des matières}%
  \renewcommand{\listfigurename}{Table des figures}%
  \renewcommand{\listtablename}{Liste des tableaux}%
  \renewcommand{\chaptername}{Chapitre}%
  \renewcommand{\appendixname}{Annexe}%
  \renewcommand{\bibname}{Références}%
  \renewcommand{\figurename}{Figure}%
  \renewcommand{\tablename}{Tableau}%
}
\hypersetup{colorlinks=true,linkcolor=blue!60!black,citecolor=blue!60!black,urlcolor=blue!60!black}
% ---- macros de notre travail
\usepackage{tabularx,array,seqsplit}
\definecolor{navy}{HTML}{0B1D3A}\definecolor{steel}{HTML}{3A6EA5}\definecolor{ice}{HTML}{E7EEF6}
\definecolor{alert}{HTML}{D62839}\definecolor{ok}{HTML}{2E8B57}\definecolor{muted}{HTML}{6B7B8C}\definecolor{amber}{HTML}{B7791F}
\newcommand{\usection}[1]{\section*{#1}\addcontentsline{toc}{section}{#1}}
\newcommand{\code}[1]{\texttt{\detokenize{#1}}}
\newcommand{\F}{\textcolor{ok}{\textbf{[F]}}}
\newcommand{\I}{\textcolor{amber}{\textbf{[I]}}}
\newcommand{\Hy}{\textcolor{alert}{\textbf{[H]}}}
\newcommand{\Pc}{\textcolor{ok}{\textbf{[C]}}}
\newcommand{\Pe}{\textcolor{amber}{\textbf{[E]}}}
\newcommand{\ie}{c.-à-d.\ }
\newcommand{\gf}{\mathbb{F}_{16}}
\newcommand{\vpshufb}{\code{vpshufb}}
\newcommand{\sha}[1]{{\ttfamily\scriptsize\seqsplit{#1}}}
\newtheorem{proposition}{Proposition}[chapter]
\newtheorem{theorem}[proposition]{Théorème}
\newtheorem{lemme}[proposition]{Lemme}
\newtheorem{corollaire}[proposition]{Corollaire}
\theoremstyle{definition}\newtheorem{definition}[proposition]{Définition}
\theoremstyle{remark}\newtheorem{remarque}[proposition]{Remarque}
\newtcolorbox{encadre}{colback=ice,colframe=ice,boxrule=0pt,arc=2pt,left=4pt,right=4pt,top=3pt,bottom=3pt,fontupper=\small}
\newcounter{lastfrontpage}


\begin{document}
\sloppy
% UoB guidlines:t
%
% Preliminary pages
% 
% The five preliminary pages must be the Title Page, Abstract, Dedication
% and Acknowledgements, Author's Declaration and Table of Contents.
% These should be single-sided.
% 
% Table of contents, list of tables and illustrative material
% 
% The table of contents must list, with page numbers, all chapters,
 % sections and subsections, the list of references, bibliography, list of
% abbreviations and appendices. The list of tables and illustrations
% should follow the table of contents, listing with page numbers the
% tables, photographs, diagrams, etc., in the order in which they appear
% in the text.
% 
%\addtocontents{toc}{\hspace{-0.60cm}{\bf D\'edicace} \nobreak \mbox{}\hfill{\bf i}\par\nobreak}
%\addtocontents{toc}{\hspace{-0.60cm}{\bf } \nobreak \mbox{}\hfill{\bf }\par\nobreak}
%\addtocontents{toc}{\hspace{-0.60cm}{\bf Remerciements} \nobreak \mbox{}\hfill{\bf ii}\par\nobreak}
%\addtocontents{toc}{\hspace{-0.60cm}{\bf } \nobreak \mbox{}\hfill{\bf }\par\nobreak}
%\addtocontents{toc}{\hspace{-0.60cm}{\bf R\'esum\'e} \nobreak \mbox{}\hfill{\bf vi}\par\nobreak}
%\addtocontents{toc}{\hspace{-0.60cm}{\bf } \nobreak \mbox{}\hfill{\bf }\par\nobreak}
%\addtocontents{toc}{\hspace{-0.60cm}{\bf Abstract} \nobreak \mbox{}\hfill{\bf vii}\par\nobreak}
%\addtocontents{toc}{\hspace{-0.60cm}{\bf } \nobreak \mbox{}\hfill{\bf }\par\nobreak}
%\addtocontents{toc}{\hspace{-0.60cm}{\bf Abr\'eviations} \nobreak \mbox{}\hfill{\bf iii}\par\nobreak}
%\addtocontents{toc}{\hspace{-0.60cm}{\bf } \nobreak \mbox{}\hfill{\bf }\par\nobreak}
%\addtocontents{toc}{\hspace{-0.60cm}{\bf Glossaire} \nobreak \mbox{}\hfill{\bf v}\par\nobreak}
%\addtocontents{toc}{\hspace{-0.60cm}{\bf } \nobreak \mbox{}\hfill{\bf }\par\nobreak}
%\addtocontents{toc}{\hspace{-0.60cm}{\bf Liste des Tableaux} \nobreak \mbox{}\hfill{\bf viii}\par\nobreak}
%\addtocontents{toc}{\hspace{-0.60cm}{\bf } \nobreak \mbox{}\hfill{\bf }\par\nobreak}
%\addtocontents{toc}{\hspace{-0.60cm}{\bf Table des Figures} \nobreak \mbox{}\hfill{\bf ix}\par\nobreak}
%\addtocontents{toc}{\hspace{-0.60cm}{\bf } \nobreak \mbox{}\hfill{\bf }\par\nobreak}
%\addtocontents{toc}{\hspace{-0.60cm}{\bf Table des matières} \nobreak \mbox{}\hfill{\bf xi}\par\nobreak}




\frontmatter
\pagenumbering{roman}
%
%
% File: Title.tex
% Author: V?ctor Bre?a-Medina
% Description: Contains the title page
%
% UoB guidelines:
% 
% At the top of the title page, within the margins, the dissertation should give the title and, if 
% necessary, sub-title and volume number. If the dissertation is in a language other than English, the 
% title must be given in that language and in English. The full name of the author should be in the 
% centre of the page. At the bottom centre should be the words ?A dissertation submitted to the 
% University of Bristol in accordance with the requirements for award of the degree of ? in the 
% Faculty of ...?, with the name of the school and month and year of submission. The word count of 
% the dissertation (text only) should be entered at the bottom right-hand side of the page.
%
%
%\hspace{5cm}
%\pagecolor{fondpaille}%\afterpage{\nopagecolor}
\bord
\begin{titlingpage}
\begin{SingleSpace}

%\calccentering{\unitlength} 
\begin{adjustwidth*}{\unitlength}{-\unitlength}
\vspace*{13mm}
\vspace{-4cm}
\centering
\begin{center}
 	\begin{tabular}{ccc}
				{\large UNIVERSITE DE YAOUNDE I} & \multirow{8}{*}{\includegraphics[ width =3cm , height =3cm]{logos/uy1.png}} & \large{UNIVERSITY OF YAOUNDE I}\\
				%& & \\
				\textbf{*******} &    & \textbf{*******} \\
				%& & \\
				\textcolor{black}{ECOLE NATIONALE SUPERIEURE}& & \textcolor{black}{NATIONAL ADVANCED SCHOOL}\\ 
				\textcolor{black}{POLYTECHNIQUE DE YAOUNDE}	  & & \textcolor{black}{OF ENGINEERING OF YAOUNDE} \\
				%& & \\
				\textbf{********}    & & \textbf{*******} \\
				%& & \\
				DEPARTEMENT DE GENIE& & DEPARTMENT OF COMPUTER \\ 
				INFORMATIQUE & & ENGINEERING \\
				\end{tabular}
	
\end{center}
\hspace{3cm}
\vspace{0.3cm}
\begin{center}
\rule[0.5ex]{\linewidth}{2pt}\vspace*{-\baselineskip}\vspace*{3.2pt}
\rule[0.6ex]{\linewidth}{1pt}\\[\baselineskip]
{\Large {\textcolor{blue}{\textsc{ \textbf{IMPLÉMENTATION ET ÉVALUATION PRATIQUE
D'UN SCHÉMA DE SIGNATURE POST-QUANTIQUE
} }}
} }\\
\rule[0.5ex]{\linewidth}{1pt}\vspace*{-\baselineskip}\vspace{3.2pt}
\rule[0.5ex]{\linewidth}{2pt}\\
\vspace{4mm}
{\Large \textbf{Mémoire de Master II}}\\
\vspace{4mm}
{\Large Présenté et soutenu par : } \\
\vspace{4mm}
{\large \textsc{\textbf{\textcolor{blue}{ANGOAH NGOAH MARCEL ERIC}}}}\\
\vspace{4mm}
{\Large En vue de l'obtention du :} \\
\vspace{4mm}
{\large \textbf{\textcolor{blue}{Master II en Génie Informatique}}}\\
\vspace{4mm}
{\Large Sous la supervision de :}\\
\vspace{4mm}

{\normalsize \textsc{\textbf{\textcolor{blue}{Hervé Talé Kalachi, Chargé de Cours, UY1}}}}\
\vspace{4mm}

\vspace{6mm}

\vspace{3mm}


\begin{tabular}{c}
{\Large \textcolor{blue}{Année académique 2025-2026}}\\ \\

\end{tabular}

%Front page corner design...read ornaments documentation for more info
%Peterson Yuhala
%\begin{tikzpicture}[remember picture, overlay, start chain, node distance=-2mm,color=orange]
%\node (nworn) [shift={(5mm,-5mm)}, anchor=north west, on chain ] at (current page.north west) {\pgfornament[width=8mm]{24}};
%\foreach \i in {1,...,21}
%\node [on chain] {\pgfornament[width=8mm]{19}};
%\node (neorn) [on chain] {\pgfornament[width=8mm]{24}};
%\foreach \i in {1,...,33}
%\node [continue chain=going below, on chain] {\pgfornament[width=8mm]{24}};
%\node (seorn) [on chain] {\pgfornament[width=8mm]{24}};
%\foreach \i in {1,...,21}
%\node [continue chain=going left, on chain] {\pgfornament[width=8mm]{24}};
%\node (sworn) [on chain] {\pgfornament[width=8mm]{24}};
%\foreach \i in {1,...,33}
%\node [continue chain=going above, on chain] {\pgfornament[width=8mm]{24}};
%\end{tikzpicture}
% \begin{tikzpicture}[remember picture, overlay]
%  \begin{scope}[shift={(current page.south west)},shift={(1,1)},scale=1]
%  \shade[ball color=blue,opacity=.6] (0,0) circle (10ex);
%  \shade[ball color=blue,opacity=.8] (1.7,1) circle (6ex);
%  \shade[ball color=blue,opacity=.8] (1.5,3) circle (2ex);
%  \shade[ball color=blue,opacity=.5] (-0.5,3) circle (1ex);
%  \shade[ball color=blue,opacity=.8] (1,4) circle (1ex);
%  \shade[ball color=blue,opacity=.6] (3.5,2.5) circle (2ex);
%  \shade[ball color=blue,opacity=.8] (2.5,3.5) circle (2ex);
%  \end{scope}
%  \end{tikzpicture}
%%%%%%%%%%%%%%%%%%%%%%%%%%%%%%%%%%%%%%%%%%%%%%%%%%%%%%%%%%%%

%\cornerplus{41}{88}
\end{center}


\end{adjustwidth*}

\end{SingleSpace}
\end{titlingpage}

\cleardoublepage

\begin{titlingpage}
\begin{SingleSpace} 

\begin{center}

% Logos
\begin{figure}[h]
    \begin{minipage}{0.4\textwidth}
        \centering
        \includegraphics[scale=0.5]{logos/enpy.png}
    \end{minipage}
    \hspace{3.5cm}
    \begin{minipage}{0.4\textwidth}
        \centering
        \includegraphics[scale=0.5]{logos/uy1.png}
    \end{minipage}
\end{figure}

% Titre du mémoire
\rule[0.5ex]{\linewidth}{2pt}\\[\baselineskip]
{\Large{\textcolor{blue}{\textsc{IMPLÉMENTATION ET ÉVALUATION PRATIQUE
D'UN SCHÉMA DE SIGNATURE POST-QUANTIQUE}}}}\\
\rule[0.5ex]{\linewidth}{2pt}\\

\vspace{6mm}
{\Large \textbf{ Mémoire de Master II}}\\
\vspace{6mm}
{\Large Présenté et soutenu par :} \\
\vspace{6mm}
{\Large \textsc{\textbf{\textcolor{blue}{ANGOAH NGOAH MARCEL ERIC}}}}\\
\vspace{8mm}
{\Large En vue de l'obtention du :} \\
\vspace{4mm}
{\large \textbf{\textcolor{blue}{Master II en Génie Informatique}}}\\

\vspace{6mm}
\begin{tabular}{c}
    {\Large \textcolor{blue}{Année académique 2025-2026}}\\ \\
    
    {\Large \textcolor{blue}{Soutenu le [date]}}
\end{tabular}

\end{center}

\end{SingleSpace}
\end{titlingpage}

\cleardoublepage
%
% file: dedication.tex
% author: V?ctor Bre?a-Medina
% description: Contains the text for thesis dedication
%

\chapter*{\Large{Dédicace}}
\label{chap:dedicace}
\addstarredchapter{Dédicace}
\markboth{Dédicace}{}
%\begin{SingleSpace}
\vspace{3cm}
\begin{center}
  \begin{tikzpicture}[color=blue,every node/.style={inner sep=0pt}]   
\node[text  width=8.5cm,align=center](Text){%
\huge{(à compléter)}
} ;
\node[shift={(-1cm,1cm)},anchor=north west](CNW)  at (Text.north west)
               {\pgfornament[width=2cm]{61}};
\node[shift={(1cm,1cm)},anchor=north east](CNE)   at (Text.north east)
               {\pgfornament[width=2cm,symmetry=v]{61}}; 
\node[shift={(-1cm,-1cm)},anchor=south west](CSW) at (Text.south west)
               {\pgfornament[width=2cm,symmetry=h]{61}}; 
\node[shift={(1cm,-1cm)},anchor=south east](CSE)  at (Text.south east)   
               {\pgfornament[width=2cm,symmetry=c]{61}};  
%\pgfornamenthline{CNW}{CNE}{north}{87}
%\pgfornamenthline{CSW}{CSE}{south}{87}
%\pgfornamentvline{CNW}{CSW}{west}{87}
%\pgfornamentvline{CNE}{CSE}{east}{87} 
\end{tikzpicture}

\end{center}
%\end{SingleSpace}
\clearpage

\clearemptydoublepage
%
\chapter*{Remerciements}
\label{chap:remerciements}
\addstarredchapter{Remerciements}
\markboth{Remerciements}{}
\begin{SingleSpace}
\begin{flushright}\itshape
« Une seule main ne peut nouer un fagot de bois. »\\
\upshape\small --- Adage africain
\end{flushright}
\vspace{0.5em}
Fort de cet adage, je prends la plume pour remercier celles et ceux dont le soutien m'a permis de mener à terme la rédaction de ce travail. Ce travail n'aurait pas été possible sans l'accompagnement de plusieurs personnes, à qui j'exprime ma sincère gratitude :
\begin{itemize}[label=\ding{42}, font=\large \color{darkblue}]
\item \textbf{Dr Hervé Talé Kalachi}, mon encadrant, pour ses relectures exigeantes et les consignes méthodologiques qui ont structuré ce travail ;
\item \emph{(à compléter : membres du jury, enseignants du Département de Génie Informatique, collègues, famille, amis.)}
\end{itemize}


\end{SingleSpace}
\clearpage

\clearemptydoublepage
%
\chapter*{Abstract}
\label{chap:abstract}
\addstarredchapter{Abstract}
\markboth{Abstract}{}
\initial{M}AYO is an Oil-and-Vinegar multivariate signature scheme submitted to post-quantum standardisation. In the official AVX2 C implementation, products of the public matrix $P_1$ by the vectors of $V$ are computed with \code{vpshufb} lookups on packed $\mathbb{F}_{16}$ vectors. When the vector length $m$ is not a multiple of 64, the official kernel covers the tail with an overlapping 32-byte load that duplicates work. We study a local change, called A2, which handles the 8-byte tail of each $P_1$ entry in a single 128-bit register and one xmm \code{vpshufb} on the low lane of the multiplication table, reducing the number of shuffles per scalar pair from 4 to 3 (MAYO-1) and from 6 to 5 (MAYO-5). We formalise the counting argument and its limits, establish the properties on which correctness rests, and measure the effect on a physical laptop with a paired protocol. Kernel gains are 12--29\,\%; the full-signature gain is 10.3\,\% for MAYO-1 but only 1.3\,\% for MAYO-5. Static analysis shows that in the signature context GCC~16.2.0 emits extra 32-byte stores in several loops, asymmetrically between the two versions, consistently with the measured loss; causality is not established. GCC's dumps show that the extra stores are per-iteration stores of the accumulator arrays that are not moved out of the loop; SRA is ruled out, and why the compiler makes this choice remains open. A targeted comparison, with the two versions compiled separately and without probes (401 process pairs), gives 10.1\,\% for MAYO-1 and 1.5\,\% for MAYO-5. The compiler's analysis files (passes \code{lim2} and \code{pre}) show that the stores leave the A2 loops in the all-probes construction and not in the probe-free one; the cause of this difference remains open.

\paragraph{Keywords: MAYO, multivariate signatures, post-quantum cryptography, AVX2, vpshufb, finite field F16, optimisation, constant time.}
\clearpage

\clearemptydoublepage
%
\chapter*{Résumé}
\label{chap:resume}
\addstarredchapter{Résumé}
\markboth{Résumé}{}
\initial{M}AYO est un schéma de signature numérique « multivarié » de type Oil-and-Vinegar. Il est candidat à la standardisation de la cryptographie post-quantique. Dans le code C officiel (MAYO-C) pour AVX2, le programme multiplie la matrice publique $P_1$ par les vecteurs de $V$ avec l'instruction \code{vpshufb}. Les vecteurs de $\mathbb{F}_{16}$ y sont rangés à deux éléments par octet. Quand la longueur $m$ n'est pas un multiple de 64, le code officiel traite la fin du vecteur en relisant 32 octets déjà traités. Cela fait des calculs en double : 48 éléments sur 128 pour MAYO-1 ($m=80$) et 50 sur 192 pour MAYO-5 ($m=142$).

Notre travail étudie une petite modification, appelée A2. Elle traite la fin de 8 octets de chaque entrée de $P_1$ dans un seul registre de 128 bits : un seul \code{vpshufb} xmm remplace deux \code{vpshufb} ymm. Le nombre de \code{vpshufb} par paire de nombres passe de 4 à 3 (MAYO-1) et de 6 à 5 (MAYO-5). Nous démontrons cette limite de comptage, nous établissons les propriétés qui rendent A2 correct, puis nous mesurons l'effet sur un ordinateur portable (Intel Core i9-9880H, GCC~16.2.0). Les mesures comparent les deux versions dans les mêmes conditions, et l'unité statistique est le processus (une exécution du programme).

Voici les résultats établis. Sur les noyaux pris seuls, A2 réduit le temps de 19 à 22\,\% (\code{P1_times_Vt}) et de 12 à 29\,\% (\code{P1t_times_Vt}). Sur la signature complète, le gain est de 10,3\,\% pour MAYO-1, mais de seulement 1,3\,\% pour MAYO-5. Pour MAYO-5, le résultat principal est donc 1,3\,\% ; les 7,2\,\% obtenus avec le binaire instrumenté (toutes les sondes actives) servent seulement à expliquer pourquoi le protocole de mesure compte. En lisant le code produit par GCC~16.2.0 (sans l'exécuter), nous voyons que, sans toutes les sondes, le compilateur ajoute dans plusieurs boucles une écriture mémoire par accumulateur (32 octets pour un registre ymm, 16 octets pour un xmm), et que ce n'est pas la même chose dans les deux versions. Les temps mesurés varient de la même façon, mais le lien de cause à effet n'est pas démontré, et un modèle à un seul paramètre exagère le gain d'environ 40\,\%. Ces écritures en plus sont les stockages des tableaux d'accumulateurs, refaits à chaque tour de boucle au lieu d'être faits une seule fois après la boucle. Nous écartons la passe SRA du compilateur comme cause. Pourquoi le compilateur fait ce choix reste une question ouverte. Une comparaison ciblée, avec les deux versions compilées séparément et sans sondes (401 paires de processus), donne 10,1\,\% pour MAYO-1 et 1,5\,\% pour MAYO-5. Les fichiers d'analyse du compilateur (passes \code{lim2} et \code{pre}) montrent que les écritures sortent des boucles de la version A2 dans la construction avec toutes les sondes, et pas dans celle sans sondes ; la cause de cette différence reste ouverte.

\paragraph{Mots-clés : MAYO, signatures multivariées, cryptographie post-quantique, AVX2, \code{vpshufb}, $\mathbb{F}_{16}$, optimisation, mesure de performance.}
\clearpage

\clearemptydoublepage
%
\chapter*{Liste des abréviations}
\label{chap:abbreviations}
\addstarredchapter{Liste des abréviations}
\markboth{Liste des abréviations}{}
\begin{itemize}
\item[•] \textbf{AVX2 :} \emph{Advanced Vector Extensions 2} : instructions des processeurs qui traitent 256 bits (32 octets) à la fois
\item[•] \textbf{CT :} \emph{Constant Time} (temps constant)
\item[•] \textbf{DSE :} \emph{Dead Store Elimination} (passe de GCC)
\item[•] \textbf{FRE :} \emph{Full Redundancy Elimination} (passe de GCC)
\item[•] \textbf{GCC :} \emph{GNU Compiler Collection}
\item[•] \textbf{GFNI :} \emph{Galois Field New Instructions}
\item[•] \textbf{IC95 :} intervalle de confiance à 95\,\%
\item[•] \textbf{KAT :} \emph{Known-Answer Test} : test de « réponses connues » (on compare à des valeurs officielles)
\item[•] \textbf{LIM :} \emph{Loop Invariant Motion} (passe de GCC)
\item[•] \textbf{MAYO :} schéma de signature multivariée de type Oil-and-Vinegar
\item[•] \textbf{FIPS :} \emph{Federal Information Processing Standards} : normes fédérales américaines de traitement de l'information
\item[•] \textbf{MQ :} \emph{Multivariate Quadratic} : équations de degré 2 à plusieurs variables
\item[•] \textbf{NIST :} \emph{National Institute of Standards and Technology} : Institut national américain des normes et de la technologie
\item[•] \textbf{OV, UOV :} \emph{(Unbalanced) Oil and Vinegar}
\item[•] \textbf{PRE :} \emph{Partial Redundancy Elimination} (passe de GCC)
\item[•] \textbf{SIMD :} \emph{Single Instruction, Multiple Data}
\item[•] \textbf{SRA :} \emph{Scalar Replacement of Aggregates} (passe de GCC)
\item[•] \textbf{TSC :} \emph{Time Stamp Counter} : compteur de ticks du processeur
\end{itemize}
\clearpage

\clearemptydoublepage

\chapter*{Glossaire}
\label{chap:Glossaire}
\addstarredchapter{Glossaire}
\markboth{Glossaire}{}
\begin{itemize}
\item[•] \textbf{$m$-vec :} vecteur de $m$ éléments de $\mathbb{F}_{16}$, empaqueté à deux éléments par octet.

\item[•] \textbf{Limbe :} mot de 64 bits d'un accumulateur.

\item[•] \textbf{Voie (\emph{lane}) :} moitié de 128 bits d'un registre ymm.

\item[•] \textbf{Table de multiplication (multab) :} table de 16 octets donnant, pour une paire de nombres, les deux produits par un quartet.

\item[•] \textbf{Corps fini :} ensemble fini de nombres où l'on peut additionner, soustraire, multiplier et diviser (par exemple $\mathbb{F}_{16}$).\item \textbf{Forme quadratique :} expression de degré 2 en plusieurs variables, par exemple $x_1x_2+x_3^2$.\item \textbf{Empreinte (hachage) :} courte suite de nombres calculée à partir d'un message ; deux messages différents ont des empreintes différentes.\item \textbf{Sel :} valeur aléatoire ajoutée au message avant le calcul de l'empreinte.\item \textbf{Compilateur, assembleur :} le compilateur transforme le code C en assembleur, la suite d'instructions que le processeur exécute.\item \textbf{Passe (du compilateur) :} une étape d'optimisation du code (par exemple SRA, PRE, DSE).\item \textbf{Bootstrap :} méthode qui rejoue le calcul sur des tirages au hasard des mesures pour estimer son incertitude.\item \textbf{Fin de vecteur :} fin d'un $m$-vec qui ne remplit pas un registre complet de 32 octets.

\item[•] \textbf{Recouvrement :} chargement reculé du dernier registre, qui retraite des octets déjà traités.

\item[•] \textbf{SAFE :} variante d'A2 qui ne lit rien au-delà de l'entrée de $P_1$.

\item[•] \textbf{Tick TSC :} unité du compteur d'horodatage (\code{rdtscp}) ; ce n'est pas un cycle cœur.

\item[•] \textbf{Processus :} unité statistique : une exécution du banc, 3\,001 paires.

\item[•] \textbf{Doublon :} écritures de 32 octets supplémentaires par accumulateur dans le code généré, selon la construction.

\item[•] \textbf{E1, E1b :} analyses statiques du code généré (E1b : origine du doublon, exécutée le 25 octobre 2026 sous GCC~16.2.0).

\item[•] \textbf{Fait, inférence, hypothèse :} statuts des énoncés, notés \F{}, \I{}, \Hy{}.

\end{itemize}
\clearpage

\clearemptydoublepage

\renewcommand{\contentsname}{Table des matières}
\maxtocdepth{subsection}

\dominitoc
\tableofcontents
%\addtocontents{toc}{\par\nobreak \mbox{}\hfill{\bf Page}\par\nobreak}
\clearemptydoublepage


\selectlanguage{french}\renewcommand{\listfigurename}{Liste des figures}\listoffigures
%\addtocontents{lof}{\par\nobreak\textbf{{\scshape Figure} \hfill Page}\par\nobreak}
\cleardoublepage
%


\renewcommand{\listtablename}{Liste des tableaux}\listoftables

%\addtocontents{lot}{\par\nobreak\textbf{{\scshape Tableau} \hfill Page}\par\nobreak}
\clearemptydoublepage
%
%
% The bulk of the document is delegated to these chapter files in
% subdirectories.
\clearemptydoublepage
\setcounter{lastfrontpage}{\value{page}}   % next page after xiii, i.e. xiv
\mainmatter
\mainmatter
%
\clearemptydoublepage
\import{chapters/intro/}{intro.tex}
\clearemptydoublepage
\import{chapters/chapter01/}{chap01.tex}
\clearemptydoublepage
\import{chapters/chapter02/}{chap02.tex}
\clearemptydoublepage
\import{chapters/chapter03/}{chap03.tex} 
\clearemptydoublepage
\import{chapters/conclusion/}{conclusion.tex}
\clearemptydoublepage

\clearemptydoublepage
\markboth{Références}{}
\begin{thebibliography}{99}
\bibitem{mayospec} MAYO Team, \emph{MAYO : spécification du troisième tour} (algorithmes~7 et~8), \url{https://pqmayo.org/assets/specs/mayo-round3-white-back.pdf}, consulté en novembre 2026.
\bibitem{lidl} R. Lidl, H. Niederreiter, \emph{Finite Fields}, 2\textsuperscript{e} éd., Cambridge University Press, 1997.
\bibitem{horn} R. A. Horn, C. R. Johnson, \emph{Matrix Analysis}, 2\textsuperscript{e} éd., Cambridge University Press, 2013.
\bibitem{lang} S. Lang, \emph{Algebra}, 3\textsuperscript{e} éd., Graduate Texts in Mathematics 211, Springer, 2002.
\bibitem{hac} A. J. Menezes, P. C. van Oorschot, S. A. Vanstone, \emph{Handbook of Applied Cryptography}, CRC Press, 1996.
\bibitem{shor} P. W. Shor, \emph{Algorithms for quantum computation: discrete logarithms and factoring}, FOCS 1994.
\bibitem{patarin} J. Patarin, \emph{The oil and vinegar signature scheme}, Dagstuhl Workshop on Cryptography, 1997.
\bibitem{kipnisshamir} A. Kipnis, A. Shamir, \emph{Cryptanalysis of the Oil and Vinegar Signature Scheme}, CRYPTO 1998.
\bibitem{uov} A. Kipnis, J. Patarin, L. Goubin, \emph{Unbalanced Oil and Vinegar Signature Schemes}, EUROCRYPT 1999.
\bibitem{beullensrainbow} W. Beullens, \emph{Breaking Rainbow Takes a Weekend on a Laptop}, CRYPTO 2022.
\bibitem{mayo} W. Beullens, \emph{MAYO: Practical Post-Quantum Signatures from Oil-and-Vinegar Maps}, IACR ePrint 2021/1144 (SAC 2021).
\bibitem{nibbling} W. Beullens, F. Campos, S. Celi, B. Hess, M. J. Kannwischer, \emph{Nibbling MAYO: Optimized Implementations for AVX2 and Cortex-M4}, IACR ePrint 2023/1683 (version consultée).
\bibitem{nistfips} NIST, \emph{FIPS 203, 204, 205: normes de chiffrement et de signature post-quantiques}, août 2024.
\bibitem{nistadd} NIST, \emph{Call for Additional Digital Signature Schemes for the Post-Quantum Cryptography Standardization Process}, 2022.
\bibitem{kocher} P. Kocher, \emph{Timing Attacks on Implementations of Diffie-Hellman, RSA, DSS, and Other Systems}, CRYPTO 1996.
\bibitem{bernstein} D. J. Bernstein, \emph{Cache-timing attacks on AES}, 2005.
\bibitem{russes} V. Arlazarov, E. Dinic, M. Kronrod, I. Faradzev, \emph{On economical construction of the transitive closure of an oriented graph}, Soviet Math.\ Doklady 11, 1970.
\bibitem{albrecht2011} M. R. Albrecht, G. V. Bard, C. Pernet, \emph{Efficient Dense Gaussian Elimination over the Finite Field with Two Elements}, arXiv:1111.6549, 2011.
\bibitem{bch23} W. Beullens, M.-S. Chen, S.-H. Hung, M. J. Kannwischer, B.-Y. Peng, C.-J. Shih, B.-Y. Yang, \emph{Oil and Vinegar: modern parameters and implementations}, IACR TCHES 2023(3).
\bibitem{mayoc} Dépôt officiel \code{PQCMayo/MAYO-C}, commit \code{74de637a7c4a7f9fafa888183bc05f5cf920cd2d}.
\end{thebibliography}\makeatother

\appendix
\clearpage
\import{chapters/appendices/}{app0A.tex}
\end{document}
